%!TEX root = document.tex

\section{软件基本信息}

\subsection{软件名称和版本信息}

\begin{table}[h]
\centering
\caption{md-highlighter 软件基本信息}
\begin{tabular}{p{3cm}p{10cm}}
\toprule
\textbf{项目} & \textbf{描述} \\
\midrule
软件全称 & 分子动力学文件语法高亮 VScode 插件 \\
软件简称 & md-highlighter \\
版本号 & 0.0.7 \\
软件分类 & 应用软件 \\
软件类型 & Visual Studio Code 扩展插件 \\
开发语言 & JSON 配置文件 \\
运行环境 & Visual Studio Code 1.70.0 及以上版本 \\
开发者 & 高旭帆（gxf1212） \\
联系方式 & gxf1212@zju.edu.cn \\
开发单位 & 个人开发者 \\
仓库地址 & \href{https://github.com/gxf1212/md-highlighter}{https://github.com/gxf1212/md-highlighter} \\
程序规模 & package.json、language-configuration.json 以及 23 个语法定义文件，共 1761 行 \\
当前公开版本 & 0.0.7 \\
\bottomrule
\end{tabular}
\end{table}

\subsection{软件用途}

md-highlighter 用于在 Visual Studio Code 中为分子动力学相关文本文件提供语法高亮。这个项目当前做的事情比较明确，不涉及计算、建模或结果分析，主要是在编辑器里完成文件类型识别、语法元素匹配和颜色作用域分配，方便直接阅读和修改专业输入文件。

从仓库实现看，它的工作流程也比较直白：先在 \lstinline|package.json| 里注册语言标识、文件扩展名和语法文件路径，再由 \lstinline|language-configuration.json| 提供通用的括号、注释和自动补全配置；各类具体格式的匹配规则写在 \lstinline|syntaxes/*.tmLanguage.json| 中，最后由 VS Code 根据匹配结果给文本片段赋予作用域，再交给当前主题显示颜色。

\subsection{适用对象}

本工具面向需要直接编辑分子模拟输入文件的用户。比较常见的使用者，一类是使用 NAMD、CHARMM、Amber、AmberTools、Gromacs、PLUMED 的研究人员；一类是在 VS Code 中查看力场、拓扑、坐标和参数文件的实验人员；另外也包括需要补充或维护这类文本语法高亮规则的开发者。

\subsection{主要功能}

当前版本和源码能直接对应上的功能，主要是这几块。第一是文件识别，项目里注册了 23 个语言 ID，并把常见扩展名或文件名模式映射到对应语言。第二是语法高亮，不同文件格式各自定义匹配规则，对关键字、原子名、残基名、编号、坐标和参数值等内容着色。第三是作用域映射，尽量复用 VS Code 主题可识别的标准 TextMate 作用域，例如 \lstinline|entity.name|、\lstinline|constant.numeric|、\lstinline|support.type|。另外，这个扩展本身比较轻，不包含后台服务、数据库或独立执行程序，安装后由编辑器直接加载。

\subsection{支持范围}

按照 \lstinline|package.json| 当前注册结果，本软件共提供 23 个语法定义，覆盖以下文件类型：

\begin{table}[!h]
\centering
\caption{已注册的语言和文件类型}
\begin{tabular}{p{2cm}p{11cm}}
\toprule
\textbf{类别} & \textbf{语言 ID 或扩展名} \\
\midrule
NAMD或CHARMM & rtf（\lstinline|.rtf|、\lstinline|.str|、\lstinline|.inp|、\lstinline|.prm|、\lstinline|.par|）、pdb（\lstinline|.pdb|）、psf（\lstinline|.psf|） \\
Amber或AmberTools & in（\lstinline|.in|，并匹配 \lstinline|leaprc.*|）、prmtop（\lstinline|.prmtop|、\lstinline|.parm7|）、inpcrd（\lstinline|.inpcrd|、\lstinline|.rst7|）、prepin（\lstinline|.prepin|、\lstinline|.prepi|）、frcmod（\lstinline|.frcmod| 及若干参数文件名模式）、lib（\lstinline|.lib|、\lstinline|.off|）、ac（\lstinline|.ac|）、mc（\lstinline|.mc|） \\
Gromacs & gro、atp、arn、rtp、hdb、r2b、tdb、mtp、vsd \\
小分子 & mol2、sdf（\lstinline|.sdf|、\lstinline|.mol|） \\
其他 & plumed（\lstinline|.plumed.dat|） \\
\bottomrule
\end{tabular}
\end{table}

\subsection{实现特点}

本软件的实现特点可以直接从代码文件看出来。现在的组织方式基本是一种语言对应一个 \lstinline|.tmLanguage.json| 文件，后续维护时比较容易单独调整。各语法文件内部主要使用内联正则表达式描述文本结构，没有额外的编译步骤。高亮范围也是基于文本模式匹配来实现，不依赖外部命令，也不读取网络数据。换句话说，当前版本主要解决的是“阅读和编辑时更容易分辨字段”这个问题，不负责语义计算、参数求解或自动修复。

\subsection{文档编写目的}

本说明文档用于说明 md-highlighter 的实际功能、安装方式、支持文件和使用边界，供软件著作权登记时与申请表、源程序文档进行对应核对。文中描述以当前仓库代码和随附截图为依据，能写到哪一步就写到哪一步，不扩展到程序里还没有实现的功能。
